\documentclass[a4paper]{article}

\usepackage{graphicx}
\usepackage{amsmath,amssymb}
\usepackage{minted}
\usepackage{multirow}
\usepackage{float}
\usepackage{todonotes}
\usepackage{forest}
\usepackage{xstring}
\usepackage{xcolor}
\usepackage[most]{tcolorbox}
\usepackage{xparse}
\usepackage{natbib}
\usepackage{algorithm}
\usepackage{algpseudocode}

\usetikzlibrary{graphs,graphdrawing}
\usegdlibrary{layered}

% for external graphviz
\input{/home/donkarlo/Dropbox/repo/nd_ai_project/src/nd_ai/knoweldge_representation/language/formal/computer/programming/tex/latex/code_clips/external_graphviz.tex}

% for tags
\input{/home/donkarlo/Dropbox/repo/nd_ai_project/src/nd_ai/knoweldge_representation/language/formal/computer/programming/tex/latex/parts/control_sequence/kind/semtag/semtag.tex}

% for links
\input{/home/donkarlo/Dropbox/repo/nd_ai_project/src/nd_ai/knoweldge_representation/language/formal/computer/programming/tex/latex/code_clips/seehere.tex}

\title{Superposition}
\date{}

\begin{document}
    \maketitle
    
    
    \section{A sample equation we are trying to solve}
        \[
            \begin{aligned}
                &M
                +
                A\cos\left(
                         \beta
                         +
                         2\arctan\left(
                                     \omega\tan\left(\frac{t-\alpha}{2}\right)
                    \right)
                \right)
                \\
                &=
                \\
                &M_1
                +
                A_1\cos\left(
                           \beta_1
                    +
                    2\arctan\left(
                                \omega_1\tan\left(\frac{t-\alpha_1}{2}\right)
                    \right)
                \right)
                \\
                &+
                M_2
                +
                A_2\cos\left(
                           \beta_2
                    +
                    2\arctan\left(
                                \omega_2\tan\left(\frac{t-\alpha_2}{2}\right)
                    \right)
                \right)
            \end{aligned}
        \]
    
    
    \section{Single FMM Approximation of a Two-Component FMM Signal}
        
        The goal is to approximate a signal generated by the sum of two FMM components using only one FMM signal. The original signal is assumed to have the following form:
        
        \[
            y(t_i)
            =
            M_0
            +
            f(t_i; A_1,\alpha_1,\beta_1,\omega_1)
            +
            f(t_i; A_2,\alpha_2,\beta_2,\omega_2)
        \]
        
        where each FMM component is defined as
        
        \[
            f(t; A,\alpha,\beta,\omega)
            =
            A
            \cos
            \left(
                \beta
                +
                2
                \arctan
                \left(
                    \omega
                    \tan
                    \left(
                        \frac{t-\alpha}{2}
                    \right)
                \right)
            \right)
        \]
        
        The approximation replaces the two-component signal with a single five-parameter FMM signal:
        
        \[
            \hat{y}(t_i; \theta)
            =
            M
            +
            f(t_i; A,\alpha,\beta,\omega)
        \]
        
        where the unknown parameter vector is
        
        \[
            \theta
            =
            \left[
                M,
                A,
                \alpha,
                \beta,
                \omega
            \right]^{T}
        \]
        
        The fitting problem is formulated as a nonlinear least-squares problem. For each sampled time point \(t_i\), the residual is the difference between the original two-component signal and the fitted single-component signal:
        
        \[
            r_i(\theta)
            =
            y(t_i)
            -
            \hat{y}(t_i; \theta)
        \]
        
        The objective function is the sum of squared residuals:
        
        \[
            J(\theta)
            =
            \sum_{i=1}^{N}
            r_i(\theta)^2
            =
            \sum_{i=1}^{N}
            \left(
                y(t_i)
                -
                \hat{y}(t_i; \theta)
            \right)^2
        \]
        
        The optimal parameter vector is obtained by minimizing this objective function:
        
        \[
            \theta^{*}
            =
            \arg\min_{\theta}
            J(\theta)
        \]
    
    
    \section{Pseudo-Code}
        
        \begin{algorithm}[htbp]
		\caption{Single FMM Approximation of a Two-Component FMM Signal}
        \begin{algorithmic}[1]
            \Require Time points \(T = \{t_1,t_2,\dots,t_N\}\)
            \Require Target signal values \(Y = \{y(t_1),y(t_2),\dots,y(t_N)\}\)
            \Require Parameter bounds for \(M\), \(A\), \(\alpha\), \(\beta\), and \(\omega\)
            \Ensure Best single FMM parameter vector \(\theta^{*}\)
            
            \State Define the single FMM model:
            \[
                \hat{y}(t_i; \theta)
                =
                M
                +
                A
                \cos
                \left(
                    \beta
                    +
                    2
                    \arctan
                    \left(
                        \omega
                        \tan
                        \left(
                            \frac{t_i-\alpha}{2}
                        \right)
                    \right)
                \right)
            \]
            
            \State Define the residual vector:
            \[
                R(\theta)
                =
                \left[
                    y(t_1)-\hat{y}(t_1;\theta),
                y(t_2)-\hat{y}(t_2;\theta),
                    \dots,
                y(t_N)-\hat{y}(t_N;\theta)
                \right]^{T}
            \]
            
            \State Define the scalar error function:
            \[
                J(\theta)
                =
                R(\theta)^{T}R(\theta)
            \]
            
            \State Use a global optimization method to find an initial parameter vector:
            \[
                \theta_{\mathrm{global}}
                =
                \arg\min_{\theta}
                J(\theta)
            \]
            
            \State Use bounded nonlinear least-squares optimization starting from \(\theta_{\mathrm{global}}\):
            \[
                \theta^{*}
                =
                \arg\min_{\theta}
                \|R(\theta)\|_2^2
            \]
            
            \State Compute the fitted signal:
            \[
                \hat{Y}
                =
                \left[
                    \hat{y}(t_1;\theta^{*}),
                    \hat{y}(t_2;\theta^{*}),
                    \dots,
                    \hat{y}(t_N;\theta^{*})
                \right]
            \]
            
            \State \Return \(\theta^{*}\), \(\hat{Y}\), and \(J(\theta^{*})\)
        \end{algorithmic}
        \end{algorithm}
    
    
    \section{Explanation of the Approximation}
        
        The original signal contains two FMM components. In general, the sum of two FMM waves is not exactly equal to a single FMM wave. Therefore, the method does not find an exact analytical replacement. Instead, it finds the best numerical approximation under a least-squares criterion.
        
        The single FMM model has only five parameters: baseline, amplitude, alpha, beta, and omega. These parameters are adjusted so that the generated single FMM curve becomes as close as possible to the target signal over all sampled time points.
        
        The approximation works by comparing the target value and the fitted value at every time point. The difference between them is called the residual. If the fitted curve is close to the target curve, the residuals are small. If the fitted curve is far from the target curve, the residuals are large.
        
        The optimization method minimizes the sum of squared residuals. Squaring the residuals has two effects. First, positive and negative errors do not cancel each other. Second, larger errors receive a stronger penalty than smaller errors.
        
        The fitting procedure has two stages. The first stage uses global optimization to search broadly through the parameter space and find a reasonable starting point. This reduces the chance of starting from a poor local solution. The second stage uses nonlinear least-squares optimization to refine the parameters locally and reduce the residuals more precisely.
        
        The final result is a single FMM wave whose shape, phase, amplitude, and baseline are chosen so that it approximates the two-component signal with minimum squared error over the sampled interval.


%        \bibliographystyle{plainnat}
%        \bibliography{/home/donkarlo/Dropbox/repo/nd_shared_research/bibliography/bibliography.bib}
\end{document}